\section{Find the Smallest Divisor Given a Threshold}
\label{sec:bs-smallest-divisor}

\problemstatement{We are given an array $\textit{nums}$ of $n$ positive
integers and an integer $\textit{threshold}$. We must choose a positive
integer \textit{divisor}, and the resulting sum
$\sum_i \lceil \textit{nums}[i] / \textit{divisor} \rceil$ must be
$\le \textit{threshold}$. Find the \textbf{smallest} such divisor.}

\begin{patternbox}
This is structurally identical to Koko Eating Bananas
(Section~\ref{sec:bs-koko}): the keywords \emph{``smallest divisor''} and
\emph{``sum of ceiling divisions must not exceed a threshold''} map
directly onto \emph{``smallest eating speed''} and \emph{``total hours
must not exceed the limit.''} Once you notice the structural identity, the
entire derivation transfers without modification -- only the variable
names change.
\end{patternbox}

\subsection*{Understanding the problem carefully}

Define $\textit{sumDivisions}(\textit{divisor}) = \sum_i \lceil
\textit{nums}[i] / \textit{divisor} \rceil$. As \textit{divisor} increases,
each term $\lceil \textit{nums}[i]/\textit{divisor} \rceil$ can only
decrease or stay the same, so $\textit{sumDivisions}$ is non-increasing in
\textit{divisor} -- the same monotonicity structure as
$\textit{hoursNeeded}$ in Section~\ref{sec:bs-koko}.

\begin{ideabox}[Candidate Space and Predicate]
The candidate space is integer divisors in $[1, \max(\textit{nums})]$
(any divisor at least as large as the largest element reduces every term
to $1$, so nothing larger is ever needed). The predicate is
$P(\textit{divisor}) = (\textit{sumDivisions}(\textit{divisor}) \le
\textit{threshold})$, and we want the smallest divisor where $P$ holds.
\end{ideabox}

\subsection*{Why the derivation transfers without modification}

There is no new reasoning required: this is a direct relabeling of
Koko Eating Bananas, with $\textit{piles} \to \textit{nums}$,
$k \to \textit{divisor}$, and $h \to \textit{threshold}$. This section
exists specifically to demonstrate that once you recognize the abstract
shape (\emph{minimize a divisor/rate subject to a sum-of-ceiling-divisions
budget}), you should reuse the exact same code template rather than
rederiving it -- pattern recognition, not just algorithmic cleverness, is
the actual skill being exercised here.

\subsection*{Algorithm}

Identical in structure to Algorithm~\ref{sec:bs-koko}'s
\textsc{MinEatingSpeed}, with the feasibility check replaced by
$\textit{sumDivisions}(\textit{mid}) \le \textit{threshold}$.

\subsection*{Dry run}

\dryrun{Take $\textit{nums} = [1, 2, 5, 9]$, $\textit{threshold} = 6$.}

\begin{itemize}
  \item $\textit{lo}=1,\textit{hi}=9$: $\textit{mid}=5$. Sum:
        $\lceil1/5\rceil+\lceil2/5\rceil+\lceil5/5\rceil+\lceil9/5\rceil
        =1+1+1+2=5 \le 6$. Feasible. $\textit{hi}=4$.
  \item $\textit{lo}=1,\textit{hi}=4$: $\textit{mid}=2$. Sum:
        $\lceil1/2\rceil+\lceil2/2\rceil+\lceil5/2\rceil+\lceil9/2\rceil
        =1+1+3+5=10 > 6$. Infeasible. $\textit{lo}=3$.
  \item $\textit{lo}=3,\textit{hi}=4$: $\textit{mid}=3$. Sum:
        $\lceil1/3\rceil+\lceil2/3\rceil+\lceil5/3\rceil+\lceil9/3\rceil
        =1+1+2+3=7 > 6$. Infeasible. $\textit{lo}=4$.
  \item $\textit{lo}=4,\textit{hi}=4$: $\textit{mid}=4$. Sum:
        $\lceil1/4\rceil+\lceil2/4\rceil+\lceil5/4\rceil+\lceil9/4\rceil
        =1+1+2+3=7 > 6$. Infeasible. $\textit{lo}=5$.
  \item $\textit{lo}=5 > \textit{hi}=4$: loop ends.
\end{itemize}

The answer is $5$.

\subsection*{C++ implementation}

\begin{lstlisting}
#include <bits/stdc++.h>
using namespace std;

long long sumDivisions(vector<int>& nums, int divisor) {
    long long sum = 0;
    for (int x : nums) {
        sum += (x + divisor - 1) / divisor; // ceil(x / divisor)
    }
    return sum;
}

int smallestDivisor(vector<int>& nums, int threshold) {
    int lo = 1, hi = *max_element(nums.begin(), nums.end());

    while (lo <= hi) {
        int mid = lo + (hi - lo) / 2;
        if (sumDivisions(nums, mid) <= threshold) {
            hi = mid - 1;
        } else {
            lo = mid + 1;
        }
    }
    return lo;
}

int main() {
    vector<int> nums = {1, 2, 5, 9};
    cout << "Smallest divisor: "
         << smallestDivisor(nums, 6) << endl;
    return 0;
}
\end{lstlisting}

\begin{complexitybox}
\textbf{Time Complexity.}
\[
  O(n \log(\max(\textit{nums}))),
\]
identical in structure to Section~\ref{sec:bs-koko}.
\textbf{Space Complexity.} $O(1)$ auxiliary space.
\end{complexitybox}

\begin{takeawaybox}
Recognizing that two differently-worded problems are the same template
underneath is just as valuable as deriving a template from scratch.
Maintain a small mental (or literal) catalogue of solved
binary-search-on-the-answer templates -- ``minimize rate subject to a
ceiling-division sum budget'' is now one entry in that catalogue, covering
both this problem and Koko Eating Bananas.
\end{takeawaybox}
